\chapter{Structural Derivatives}
\label{ch:bf-structural-derivatives}

Kalman derivative recursions assume that the derivatives of the state-space
objects are already available.  Structural derivatives are the upstream layer:
they map model parameters into
\[
  c_\theta,\ F_\theta,\ Q_\theta,\ a_\theta,\ H_\theta,\ R_\theta
\]
and their first and second derivatives.

\section{Provider Map}
\label{sec:bf-structural-provider-map}

Let $\psi\in\mathbb{R}^p$ denote the parameter vector used by the numerical
provider.  In HMC this is usually an unconstrained or transformed parameter
vector, not necessarily the raw economic vector.  The structural provider
declares the map
\begin{equation}
\label{eq:bf-structural-map}
  g(\psi)
  =
  \ve\!\left(
    c(\psi),F(\psi),Q(\psi),a(\psi),H(\psi),R(\psi),
    m_0(\psi),P_0(\psi)
  \right).
\end{equation}
The structural map contract uses the reduced-form objects
$(c,F,Q,a,H,R,m_0,P_0)$ because the filtering scan cannot be differentiated
without an explicit initial-condition policy.

For every active coordinate $\psi_i$, the first-order provider contract is
\begin{equation}
\label{eq:bf-structural-first-provider}
  \dot G^{(i)}
  =
  \left(
    \dot c^{(i)},\dot F^{(i)},\dot Q^{(i)},
    \dot a^{(i)},\dot H^{(i)},\dot R^{(i)},
    \dot m_0^{(i)},\dot P_0^{(i)}
  \right).
\end{equation}
For every active pair $(\psi_i,\psi_j)$, the second-order contract is
\begin{equation}
\label{eq:bf-structural-second-provider}
  \ddot G^{(ij)}
  =
  \left(
    \ddot c^{(ij)},\ddot F^{(ij)},\ddot Q^{(ij)},
    \ddot a^{(ij)},\ddot H^{(ij)},\ddot R^{(ij)},
    \ddot m_0^{(ij)},\ddot P_0^{(ij)}
  \right).
\end{equation}
Equations~\eqref{eq:bf-structural-first-provider}--\eqref{eq:bf-structural-second-provider}
are the objects consumed by Chapters~\ref{ch:bf-kalman-score} and
\ref{ch:bf-kalman-hessian}.  A backend may store them as dense tensors, sparse
blocks, named arrays, or lazy linear operators, but the contract is mathematical:
each block must mean a derivative of the same state-space law used to evaluate
the likelihood.

For BayesFilter, a derivative provider must declare:
\begin{itemize}
  \item parameter order and units;
  \item shape of every first and second derivative tensor;
  \item whether derivatives are analytic, autodiff, finite-difference, or
    unavailable;
  \item which parameters affect which reduced-form objects;
  \item whether cross-derivatives are implemented or intentionally zero.
\end{itemize}

\section{Parameter Transforms}
\label{sec:bf-structural-parameter-transforms}

Let $\phi=T(\psi)$ be the raw economic parameter vector and let
$M(\phi)$ be any reduced-form block among $c,F,Q,a,H,R,m_0,P_0$.  Write
\[
  M_\alpha=\frac{\partial M}{\partial\phi_\alpha},
  \qquad
  M_{\alpha\beta}
  =
  \frac{\partial^2M}{\partial\phi_\alpha\partial\phi_\beta},
\]
and
\[
  T_{\alpha,i}=\frac{\partial\phi_\alpha}{\partial\psi_i},
  \qquad
  T_{\alpha,ij}
  =
  \frac{\partial^2\phi_\alpha}{\partial\psi_i\partial\psi_j}.
\]
Then
\begin{align}
\label{eq:bf-structural-transform-first}
  \dot M^{(i)}
  &=
  \sum_\alpha M_\alpha T_{\alpha,i},\\
\label{eq:bf-structural-transform-second}
  \ddot M^{(ij)}
  &=
  \sum_{\alpha,\beta}
    M_{\alpha\beta}T_{\alpha,i}T_{\beta,j}
  +\sum_\alpha M_\alpha T_{\alpha,ij}.
\end{align}
These are the BayesFilter versions of the MacroFinance transformed-parameter
chain-rule equations.  They matter because HMC often runs on
$\psi=\log\sigma$, $\psi=\log\kappa$, Cholesky coordinates, Fisher transforms,
or stationarity transforms.  A correct derivative with respect to a raw standard
deviation is wrong for HMC if it is passed off as the derivative with respect to
the log standard deviation.

\section{Initial Conditions}
\label{sec:bf-structural-initial-conditions}

The initial moments in \eqref{eq:bf-structural-map} are part of the model law.
There are three common cases:
\begin{enumerate}[label=(\roman*)]
  \item fixed diffuse or fixed proper initialization:
    $\dot m_0=\dot P_0=\ddot m_0=\ddot P_0=0$ because the initial moments are
    data of the filtering problem;
  \item stationary initialization:
    $m_0$ and $P_0$ solve model equations and must be differentiated together
    with those equations;
  \item estimated or hierarchical initialization:
    $m_0$ and $P_0$ have their own parameter blocks and cross derivatives.
\end{enumerate}
For a stationary linear system with $m_0=c+Fm_0$, the first derivative satisfies
\begin{equation}
\label{eq:bf-structural-stationary-mean-first}
  (I-F)\dot m_0^{(i)}
  =
  \dot c^{(i)}+\dot F^{(i)}m_0.
\end{equation}
For $P_0=FP_0F^\top+Q$, the first derivative satisfies the Lyapunov derivative
equation
\begin{equation}
\label{eq:bf-structural-stationary-cov-first}
  \dot P_0^{(i)}
  =
  F\dot P_0^{(i)}F^\top
  +\dot F^{(i)}P_0F^\top
  +FP_0\dot F^{(i)\top}
  +\dot Q^{(i)}.
\end{equation}
Second derivatives follow by differentiating
\eqref{eq:bf-structural-stationary-mean-first} and
\eqref{eq:bf-structural-stationary-cov-first}.  BayesFilter should record
whether these equations were solved analytically, by a differentiable Lyapunov
solver, by autodiff through the stationary solver, or intentionally not used.

\section{Sparsity, Masks, and Zero Blocks}
\label{sec:bf-structural-sparsity-masks}

A derivative tensor entry equal to zero is different from a derivative entry
that is unavailable.  The provider must distinguish:
\[
  \ddot M^{(ij)}=0\quad\hbox{declared structural zero},
  \qquad
  \ddot M^{(ij)}\ \hbox{not implemented}.
\]
Structural zeros are useful for large models because they allow the score and
Hessian scan to skip blocks without changing the mathematical target.  Missing
entries are not a mathematical shortcut; they mean the backend cannot certify
the requested derivative object.

Observation masks are data-side selections.  If $J_t$ selects observed
coordinates at time $t$, the provider may expose either the full
$(a_t,H_t,R_t)$ together with $J_t$ or the already-masked blocks
\[
  a_t^o=J_ta_t,\qquad H_t^o=J_tH_t,\qquad R_t^o=J_tR_tJ_t^\top .
\]
The same selection applies to derivatives.  Since $J_t$ is data, not a
parameter, no derivative of the mask is taken.  This convention is the structural
upstream counterpart of the masking rule in
Section~\ref{sec:bf-score-masks}.

\section{Metadata and Shape Contract}
\label{sec:bf-structural-metadata-shapes}

The provider must make the following metadata inspectable before likelihood
evaluation:
\begin{itemize}
  \item names and order of $\psi_i$;
  \item transform type and raw unit for each parameter;
  \item state and observation coordinate names;
  \item dimensions of each reduced-form block;
  \item tensor layout conventions, including whether Hessian tensors store all
    pairs or only the symmetric upper triangle;
  \item symmetry and positive-definiteness policies for $Q$, $R$, and $P_0$;
  \item deterministic or algebraic state-coordinate roles when the provider is
    used by a structural nonlinear filter.
\end{itemize}
These declarations are not decorative.  Without them, a finite likelihood can be
computed from a different parameter ordering, unit system, or deterministic
state convention than the derivative tensors assume.

The MacroFinance large-scale derivative-provider tests check derivative shapes,
finite-difference parity by object type, reduced-form rank diagnostics, and
parameter-unit rescaling.  BayesFilter should preserve these as public contract
tests before a structural model is allowed to feed an HMC target.

\section{Provider Validation}
\label{sec:bf-structural-provider-validation}

The minimum validation ladder for a structural provider is:
\begin{enumerate}[label=(\roman*)]
  \item finite value check: all blocks in \eqref{eq:bf-structural-map} are finite,
    conformable, symmetric where required, and admissible before filtering;
  \item shape check:
    \eqref{eq:bf-structural-first-provider}--\eqref{eq:bf-structural-second-provider}
    have the declared dimensions for every active parameter and pair;
  \item transform check:
    \eqref{eq:bf-structural-transform-first}--\eqref{eq:bf-structural-transform-second}
    agree with finite differences in transformed coordinates;
  \item block finite-difference check: each reduced-form object is checked
    before it is inserted into a long filtering scan;
  \item end-to-end check: the resulting likelihood score and Hessian agree with
    finite differences or autodiff on small systems where that comparison is
    reliable;
  \item diagnostic check: structural zeros, unavailable derivatives,
    regularization, rejected parameter points, and mask policies are reported.
\end{enumerate}
Only after these checks pass should the structural provider feed an HMC target
or an observed-information diagnostic.  This is the upstream analogue of the
factor parity gates in Chapter~\ref{ch:bf-factor-derivatives}.
